% GENERATED FILE. Edit canonical lessons, then run npm run generate:books.
% canonical-source-hash: fnv1a64:504b7d83cefb59a2
% canonical-lessons: LA-C507-superbia, LA-C507-arrogans, LA-C507-arrogantia, LA-C507-insolens, LA-C507-vanitas

\chapter{Pride and arrogance}
\label{ch:la-pride-and-arrogance}
\hlchaptermodality{507}

\begin{chapteropening}
\textbf{By the end of this chapter:} \emph{I can describe proud and arrogant people.}

Complete the last lesson of chapter 507: I can describe proud and arrogant people.
\end{chapteropening}

\section[superbia, superbiae]{superbia, superbiae — pride, haughtiness}

\label{lesson:LA-C507-superbia}



\begin{quote}
Before the new one: say the Latin for country manners, then the Latin for generosity.
\end{quote}

\subsection*{You'll want to know: superbia, superbiae}
\textbf{superbia, superbiae} — \textquotedblleft{}pride, haughtiness\textquotedblright{}.

\begin{quote}
Superbia rēgis cīvēs vexat.
\end{quote}

\begin{quote}
\emph{The king's arrogance troubles the citizens.}
\end{quote}

Feminine, first declension, from \emph{superbus}. Usually a fault in Latin, not a virtue.

\begin{cousinweb}
\textbf{superbia} is built on \emph{superbus} \textquotedblleft{}proud\textquotedblright{}, which Lewis \& Short and EtymDB take from \emph{super} \textquotedblleft{}above\textquotedblright{}: one who holds himself above others.
\end{cousinweb}

\begin{grammarlens}[title={What you've built}]
One more word for telling, explaining and describing at B1.
\end{grammarlens}

\subsection*{Your turn}
\begin{itemize}
\raggedright
  \item \emph{Say it:} \emph{superbia, superbiae}
  \item \emph{Say it:} \emph{superbia, superbiae}, once more
  \item \emph{Say it:} a sentence of your own with \emph{superbia, superbiae}
  \item \emph{Recall:} say the Latin for country manners, then the Latin for generosity, then say \emph{superbia, superbiae} again
\end{itemize}

\begin{tcolorbox}[breakable,colback=teal!4,colframe=teal!35!black,title={Before you move on}]
What does superbia, superbiae mean? (\textquotedblleft{}Pride, haughtiness\textquotedblright{}.) Say it once more.
\end{tcolorbox}
\section[arrogāns, arrogantis]{arrogāns, arrogantis — arrogant, presumptuous}

\label{lesson:LA-C507-arrogans}



\begin{quote}
Before the new one: say the Latin for generosity, then the Latin for pride.
\end{quote}

\subsection*{You'll want to know: arrogāns, arrogantis}
\textbf{arrogāns, arrogantis} — \textquotedblleft{}arrogant, presumptuous\textquotedblright{}.

\begin{quote}
Iuvenis arrogāns nēminem audit.
\end{quote}

\begin{quote}
\emph{The arrogant young man listens to no one.}
\end{quote}

One ending. The present participle of \emph{arrogō} — to claim for oneself. English \emph{arrogant}.

\begin{cousinweb}
\textbf{arrogāns} is the present participle of \emph{arrogō} \textquotedblleft{}to claim for oneself\textquotedblright{}, built on \emph{ad-} and \emph{rogō} \textquotedblleft{}to ask\textquotedblright{} (Lewis \& Short). English \emph{arrogant} comes from it.
\end{cousinweb}

\begin{grammarlens}[title={What you've built}]
One more word for telling, explaining and describing at B1.
\end{grammarlens}

\subsection*{Your turn}
\begin{itemize}
\raggedright
  \item \emph{Say it:} \emph{arrogāns, arrogantis}
  \item \emph{Say it:} \emph{arrogāns, arrogantis}, once more
  \item \emph{Say it:} a sentence of your own with \emph{arrogāns, arrogantis}
  \item \emph{Recall:} say the Latin for generosity, then the Latin for pride, then say \emph{arrogāns, arrogantis} again
\end{itemize}

\begin{tcolorbox}[breakable,colback=teal!4,colframe=teal!35!black,title={Before you move on}]
What does arrogāns, arrogantis mean? (\textquotedblleft{}Arrogant, presumptuous\textquotedblright{}.) Say it once more.
\end{tcolorbox}
\section[arrogantia, arrogantiae]{arrogantia, arrogantiae — arrogance, conceit}

\label{lesson:LA-C507-arrogantia}



\begin{quote}
Before the new one: say the Latin for pride, then the Latin for arrogant.
\end{quote}

\subsection*{You'll want to know: arrogantia, arrogantiae}
\textbf{arrogantia, arrogantiae} — \textquotedblleft{}arrogance, conceit\textquotedblright{}.

\begin{quote}
Arrogantia eius omnibus molesta est.
\end{quote}

\begin{quote}
\emph{His arrogance annoys everyone.}
\end{quote}

Feminine, first declension. English \emph{arrogance}.

\begin{cousinweb}
\textbf{arrogantia} is built on \emph{arrogāns} \textquotedblleft{}arrogant\textquotedblright{} (Lewis \& Short; EtymDB). English \emph{arrogance} comes from it.
\end{cousinweb}

\begin{grammarlens}[title={What you've built}]
One more word for telling, explaining and describing at B1.
\end{grammarlens}

\subsection*{Your turn}
\begin{itemize}
\raggedright
  \item \emph{Say it:} \emph{arrogantia, arrogantiae}
  \item \emph{Say it:} \emph{arrogantia, arrogantiae}, once more
  \item \emph{Say it:} a sentence of your own with \emph{arrogantia, arrogantiae}
  \item \emph{Recall:} say the Latin for pride, then the Latin for arrogant, then say \emph{arrogantia, arrogantiae} again
\end{itemize}

\begin{tcolorbox}[breakable,colback=teal!4,colframe=teal!35!black,title={Before you move on}]
What does arrogantia, arrogantiae mean? (\textquotedblleft{}Arrogance, conceit\textquotedblright{}.) Say it once more.
\end{tcolorbox}
\section[īnsolēns, īnsolentis]{īnsolēns, īnsolentis — insolent, haughty; unaccustomed}

\label{lesson:LA-C507-insolens}



\begin{quote}
Before the new one: say the Latin for arrogant, then the Latin for arrogance.
\end{quote}

\subsection*{You'll want to know: īnsolēns, īnsolentis}
\textbf{īnsolēns, īnsolentis} — \textquotedblleft{}insolent, haughty; unaccustomed\textquotedblright{}.

\begin{quote}
Mīles īnsolēns agricolīs male dīcit.
\end{quote}

\begin{quote}
\emph{The insolent soldier speaks abusively to the farmers.}
\end{quote}

One ending. First unaccustomed (\emph{in-} + \emph{soleō}), then behaving beyond the usual bounds. English \emph{insolent}.

\begin{cousinweb}
\textbf{īnsolēns} is built on \emph{in-} \textquotedblleft{}not\textquotedblright{} and \emph{solēns}, the participle of \emph{soleō} \textquotedblleft{}to be accustomed\textquotedblright{} (Lewis \& Short). English \emph{insolent} comes from it.
\end{cousinweb}

\begin{grammarlens}[title={What you've built}]
One more word for telling, explaining and describing at B1.
\end{grammarlens}

\subsection*{Your turn}
\begin{itemize}
\raggedright
  \item \emph{Say it:} \emph{īnsolēns, īnsolentis}
  \item \emph{Say it:} \emph{īnsolēns, īnsolentis}, once more
  \item \emph{Say it:} a sentence of your own with \emph{īnsolēns, īnsolentis}
  \item \emph{Recall:} say the Latin for arrogant, then the Latin for arrogance, then say \emph{īnsolēns, īnsolentis} again
\end{itemize}

\begin{tcolorbox}[breakable,colback=teal!4,colframe=teal!35!black,title={Before you move on}]
What does īnsolēns, īnsolentis mean? (\textquotedblleft{}Insolent, haughty; unaccustomed\textquotedblright{}.) Say it once more.
\end{tcolorbox}
\section[vānitās, vānitātis]{vānitās, vānitātis — boastfulness, vanity; emptiness}

\label{lesson:LA-C507-vanitas}



\begin{quote}
Before the new one: say the Latin for arrogance, then the Latin for insolent.
\end{quote}

\subsection*{You'll want to know: vānitās, vānitātis}
\textbf{vānitās, vānitātis} — \textquotedblleft{}boastfulness, vanity; emptiness\textquotedblright{}.

\begin{quote}
Vānitās eius rīdicula est.
\end{quote}

\begin{quote}
\emph{His boastfulness is ridiculous.}
\end{quote}

Feminine, third declension, from \emph{vānus} — empty. English \emph{vanity}.

\begin{cousinweb}
\textbf{vānitās} is built on \emph{vānus} \textquotedblleft{}empty\textquotedblright{} (Lewis \& Short), which EtymDB traces to a Proto-Indo-European root meaning \textquotedblleft{}empty\textquotedblright{}. English \emph{vanity} and \emph{vain} come from the same family.
\end{cousinweb}

\begin{grammarlens}[title={What you've built}]
One more word for telling, explaining and describing at B1.
\end{grammarlens}

\subsection*{Your turn}
\begin{itemize}
\raggedright
  \item \emph{Say it:} \emph{vānitās, vānitātis}
  \item \emph{Say it:} \emph{vānitās, vānitātis}, once more
  \item \emph{Say it:} a sentence of your own with \emph{vānitās, vānitātis}
  \item \emph{Recall:} say the Latin for arrogance, then the Latin for insolent, then say \emph{vānitās, vānitātis} again
\end{itemize}

\begin{tcolorbox}[breakable,colback=teal!4,colframe=teal!35!black,title={Before you move on}]
What does vānitās, vānitātis mean? (\textquotedblleft{}Boastfulness, vanity; emptiness\textquotedblright{}.) Say it once more.
\end{tcolorbox}
